% Source: https://artofproblemsolving.com/community/c3h1555667_hard_counting_question_please_help
% Stable identity: aops-topic-1555667-q02
Use coordinates in which the north-south streets are $x=1,\dots,n$, the east-west streets are $y=1,\dots,n$, the cars start on the lines $x=0$, $x=n+1$, $y=0$, $y=n+1$, and the speed is $1$. Two cars can only meet at an intersection of their streets (each street has only one car). The figure shows the grid for $n=4$, with the diagonal $i=j$ (dashed blue) and the antidiagonal $i+j=n+1$ (dashed red).

\begin{asy}
usepackage("amsmath");
size(4.5cm);
defaultpen(fontsize(10pt));
int n=4;
for(int i=1;i<=n;++i){ draw((i,0.4)--(i,n+0.6),gray(0.6)); draw((0.4,i)--(n+0.6,i),gray(0.6)); }
for(int i=1;i<=n;++i){ label("$"+string(i)+"$",(i,0.4),S); label("$"+string(i)+"$",(0.4,i),W); }
label("$x$",(n+0.8,0.4),S);label("$y$",(0.4,n+0.8),W);
draw((0.5,0.5)--(n+0.5,n+0.5),dashed+blue); draw((0.5,n+0.5)--(n+0.5,0.5),dashed+red);
label("\mbox{}",(min(currentpicture,user=true)+max(currentpicture,user=true))/2,black); // keeps the text after the figure black when LaTeX typesets the labels inline
\end{asy}

\textbf{When do cars meet?} The car on the east-west street $y=j$ reaches the intersection $(i,j)$ at time $i$ if it starts at the west end ($W$) and at time $n+1-i$ if it starts at the east end ($E$). The car on the north-south street $x=i$ reaches $(i,j)$ at time $j$ from the south end ($S$) and $n+1-j$ from the north end ($N$). So the cars on $y=j$ and $x=i$ collide exactly when
\begin{itemize}
\item $i=j$ and the pair of ends is $(W,S)$ or $(E,N)$ (on the diagonal, dashed blue), or
\item $i+j=n+1$ and the pair of ends is $(W,N)$ or $(E,S)$ (on the antidiagonal, dashed red).
\end{itemize}

\textbf{The conditions.} Write $h_j\in\{W,E\}$ for the car of street $y=j$ and $v_i\in\{S,N\}$ for the car of street $x=i$. Avoiding diagonal collisions forces $v_j=N$ if $h_j=W$ and $v_j=S$ if $h_j=E$. Avoiding antidiagonal collisions forces $v_{n+1-j}=S$ if $h_j=W$ and $v_{n+1-j}=N$ if $h_j=E$. Applying the first rule to the street $y=n+1-j$ as well, these say exactly:
\[
v_j\text{ is determined by }h_j,\qquad\text{and}\qquad h_{n+1-j}\ne h_j\ \text{ for every }j .
\]
\begin{itemize}
\item If $n$ is odd, the middle street $j=\dfrac{n+1}{2}$ would need $h_j\ne h_j$: there is no valid placement.
\item If $n=2m$ is even, the ends $h_1,\dots,h_m$ may be chosen freely; then $h_{m+1},\dots,h_n$ and all the $v_i$ are forced, and every condition holds.
\end{itemize}
So the number of placements is
\[
\boxed{2^{\dfrac {n}{2}}\text{ if }n\text{ is even, and }0\text{ if }n\text{ is odd}.}
\]
For $n=2$ the two placements are the example in the statement and its mirror image.
